\scanpage{77}
\mspage{686}

\Needspace{12\baselineskip}
$7^\circ)$ Les groupes $\underline N_\rho$, \deleted{$\unsure{\underline N'_\rho}$} etc, $\underline N_\rho^{\unsure{\natural|\sharp}}$ etc.
\margin{$\unsure{\sigma\,\mathrm D\,\sigma}$}

Étudions le sous-groupe
\[
(152)\qquad
\begin{gathered}
\underline N_\rho\subset\underline G_{\rho,\sigma}:\\[-2pt]
\Vert\\[-2pt]
\{\underline u\in\Gamma\underline G_{\rho,\sigma}\mid u(\rho)=\rho^{\varepsilon_3(\underline u)}\}.
\end{gathered}
\]
\margin{Notation provisoire, cf plus bas p.697}

\deleted{On} L'hom. composé
\diagram{\begin{tikzcd}[ampersand replacement=\&,column sep=large]
\underline N_\rho \arrow[r,hook] \arrow[rr,bend right=15,"{\underline\chi_{\deleted{\unsure{\rho}}}}"'] \& \underline G_{\rho,\sigma} \arrow[r,"\underline\chi"] \& \mathbb O_{12}^{*}
\end{tikzcd}}
est épi, \deleted{car} vérifions-le pour nous rassurer.

Soit $\underline\mu$ section locale de $\mathbb O_{12}^{*}$, correspondant à
$\deleted{\unsure{\underline\mu|\nu_0}}\inserted{\nu_2,\nu_3,\deleted{\unsure{\nu_0}}}\unsure{\nu_0},\mu,\nu',\nu'',\nu\deleted{\unsure{,\nu_2,\nu_3}}\ldots$
On aura
\[
\deleted{(153)}\qquad
u_{\underline\mu}(\rho)=\lambda_1^\nu\rho
=(\lambda_1^{3\nu''}\lambda_1^{-\nu_3})\rho
=\operatorname{int}(\delta^{\nu''}\tau^{\nu_3})(\rho).
\]
\deleted{distinguons les deux cas,}

\deleted{\unsure{si} $\nu_3=0$ \unsure{$\Leftrightarrow$} i.e. $\nu=3\nu''$, on aura donc}
\[
\deleted{
\begin{gathered}
u_{\underline\mu}(\rho)=(\underbrace{\lambda_1}_{\unsure{(1-\rho_{\underline M})(\delta)\Leftarrow}})^{3\nu''}\rho
=\underbrace{\bigl((1-\rho_M)(\deleted{\unsure{\lambda_1}}\delta^{\nu''})\bigr)}_{\delta^{\nu''}\rho(\delta^{\nu''})^{-1}}\rho
=\operatorname{int}(\delta^{\nu''})(\rho)\\
\unsure{\text{i.e.}}\quad \delta=\lambda_1-\lambda_0
\end{gathered}}
\]
d'où
\[
(153)\qquad
\operatorname{int}\{(\delta^{\nu''}\tau_{\unsure{1}}^{\nu_3})^{-1}\cdot u_{\underline\mu}\}(\rho)=\rho
\]
i.e.
\[
(154)\qquad
\left\{
\begin{aligned}
\underbrace{(\delta^{\nu''}\tau_{\unsure{1}}^{\nu_3})^{-1}u_{\underline\mu}}_{u}
&=\tau_{\unsure{1}}^{\nu_3}\delta^{-\nu''}u_{\underline\mu}\in\Gamma\underline N_\rho\\
\underline\chi(u)&=\deleted{\unsure{\mu}}\underline\chi\bigl((\delta^{\nu''}\tau_1^{\nu_3})^{-1}\bigr)\cdot\underbrace{\underline\chi(u_{\underline\mu})}_{\underline\mu}\\
&\hspace{-1em}\underline\chi(\tau_1^{\nu_3})=\underbrace{(-1)^{\nu_3}}_{\varepsilon_3}
\end{aligned}
\right.
\]
\margin{\uline{NB.} On note $\underline\chi$ l'hom. can.
\diagram{\begin{tikzcd}[ampersand replacement=\&]
\underline G_{\rho,\sigma}\arrow[r] \& \mathbb O_{12}^{*}
\end{tikzcd}}}

Donc posant
\[
v=(\omega\tau')^{\nu_3}u
=\bigl(\underbrace{\omega\deleted{\unsure{\tau}}\tau'\tau}_{\sigma}\bigr)^{\nu_3}\delta^{-\nu''}u_{\underline\mu}\qquad\deleted{\text{i.e.}}
\]
\deleted{on aura}
\inserted{on aura}
\[
(154)\qquad
\underline\chi(v)=\underline\mu,\quad\text{avec }v=\sigma^{\nu_3}\delta^{-\nu''}u_{\underline\mu}\in\Gamma\underline N_\rho
\]
\[
\left(\delta=\lambda_1\lambda_0^{-1},\quad
\begin{gathered}\nu_3=\nu_3(\underline\mu),\\\nu''=\nu''(\underline\mu)\end{gathered}\right).
\]
\margin{76}

\scanpage{78}

On aura donc une suite exacte
\begin{samepage}
\[
(155)
\]
\diagram{\begin{tikzcd}[ampersand replacement=\&,column sep=large]
1\arrow[r,dash] \& \underline{\unsure{Z|N}}_\rho^{\deleted{\unsure{0}}}\arrow[r,dash] \& \underline N_\rho\arrow[r,dash] \& \mathbb O_{12}^{*}\arrow[r,dash]\arrow[d,equal] \& 1\\
\& \& \& \gamma_{\rho,\sigma}
\end{tikzcd}}
\end{samepage}
et
\[
\underline{\unsure{Z|N}}_\rho^{\deleted{\unsure{0}}}
=\underline{\operatorname{Centr}}_{\underline{\unsure{\mathcal G|\widetilde G}}^{+}}(\rho).
\]

On a, si
\[
(156)\qquad\left\{
\begin{aligned}
H^+_{S_0}&\overset{\mathrm{dfn}}{=}
\underline{\unsure{\mathcal G|\widetilde G}}^{+}/
\underbrace{\deleted{\unsure{\underline Z}}\underline{\unsure{\mathcal G|\widetilde G}}^{0}}_{\substack{\text{comp.}\\\text{neutre,}\\\simeq\underline M}}
\simeq(H^+)_{S_0}\\
\text{i.e.}\quad H^+&=\unsure{\mathcal G|\widetilde G}^{+}/M,
\end{aligned}\right.
\]
un hom. canonique
\begin{samepage}
\[
(157)
\]
\diagram{\begin{tikzcd}[ampersand replacement=\&,column sep=large]
\underline Z_\rho\arrow[r] \& \underline{\operatorname{Cent}}_{H^+_{S_0}}(\rho)\simeq\deleted{\unsure{\underline Z}}\underline L_\rho\quad(\simeq(\mathbb Z/6\mathbb Z)_{S_0})
\end{tikzcd}}
\end{samepage}
\deleted{qui est \unsure{surjectif}\inserted{épi} (le noyau est égal}
qui est épi (puisqu'il y a $\overline{\mathrm m}$ une section, provenant de l'inclusion $\underline L_\rho\subset\underline Z_\rho$) et dont le noyau est
\[
(158)\qquad
\underline Z_\rho^0=\underline Z_\rho\cap\underline M
=\{x\in\Gamma\underline M\mid\rho_{\underline M}\underline x=x\quad\text{i.e.}\quad(\unsure{1-\rho_{\underline M}|1+\rho_{\underline M}})x\deleted{\unsure{\rho}}=0\}.
\]
Ce noyau est nul au dessus de $\underbrace{\operatorname{Spec}\mathbb Z}_{S_0}\setminus\{3\}$,
mais de dim $1$ en dim $3$. Pourtant on sent bien que la ``bonne'' définition de $\underline Z_\rho^{\unsure{0}}$ \deleted{est} doit \deleted{être par}\inserted{correspondre à} une équation
\[
\deleted{\rho_{\underline M}\underline x}\qquad
(159)\qquad \unsure{\ell|\varepsilon}(1-\rho_{\underline M})x=0\qquad\uuline{\text{dans }\underline M_0},
\]
\margin{77}

\scanpage{79}
\mspage{687}

ce qui revient en fait \inserted{simplement,} à remplacer $\underline N_\rho$ défini dans (152), par l'adhérence schématique de sa restriction au dessus de $S_0-\{3\}$; On trouve donc un sous-groupe meilleur, que désormais je note $\underline N_\rho$, s'insérant dans la suite exacte
\begin{samepage}
\[
(159)
\]
\diagram{\begin{tikzcd}[ampersand replacement=\&,column sep=large]
1\arrow[r] \& \underline L_\rho\arrow[r,phantom,"{\unsure{\hookrightarrow|\leftarrow}}" description] \& \underline N_\rho\arrow[r,"\underline\chi_\rho"] \& \underline{\mathbb O}_{12}^{*}\arrow[r] \& 1
\end{tikzcd}}
\end{samepage}

Posons
\[
(160)\qquad\left\{
\begin{aligned}
\underline G_{\rho,\sigma}^{\unsure{\natural|\sharp}}
&=\underbrace{\underline{\mathbb O}_{12}^{*}}_{\gamma_{\rho,\sigma}}\cdot
\underbrace{\underline{\unsure{\mathcal G|\widetilde G}}^{\unsure{\natural|\sharp}}}_{\underline M}
\subset\underline{\mathbb O}_{12}^{*}\cdot\underline{\unsure{\mathcal G|\widetilde G}}^{+}
=\underline G_{\rho,\sigma},\\
\underline N_\rho^{\unsure{\natural|\sharp}}
&=\underline N_\rho\cap\underline G_{\rho,\sigma}^{\unsure{\natural|\sharp}}
\end{aligned}\right.
\]
on trouve alors
\[
(161)\qquad
\underline N_\rho^{\unsure{\natural|\sharp}}\cap\underline L_\rho
=\underline M\cap\underline L_\rho=\{1\}
\]
donc
\begin{samepage}
\[
(161)
\]
\diagram{\begin{tikzcd}[ampersand replacement=\&,column sep=large]
\underline N_\rho^{\unsure{\natural|\sharp}}\arrow[r,hook] \& \underbrace{\underline{\mathbb O}_{12}^{*}}_{\gamma_{\rho,\sigma}}
\end{tikzcd}}
\end{samepage}
Je dis que l'image est
\[
(162)\qquad
\underline\gamma'_{\rho,\sigma}
=\underline{\mathbb O}_{12}^{\prime *}
\overset{\mathrm{dfn}}{=}
\left\{\underline\mu\in\Gamma\underline{\mathbb O}_{12}^{*}\;\middle|\;
\begin{gathered}\varepsilon_3(\underline\mu)=1\\\text{i.e. }\nu_3(\underline\mu)=0\end{gathered}\right\}.
\]
\margin{donc on \unsure{a} $\underline N_\rho^{\unsure{\natural|\sharp}}\subset\underline Z_\rho$, il vaut mieux le noter $\underline Z_\rho^{\unsure{\natural|\sharp}}$.}
Que l'image contienne $\underline{\mathbb O}_{12}^{\prime *}$ résulte de (154), montrons qu'elle est contenue.

Soit
\[
u=f^{-1}u_{\underline\mu},\qquad
f\in\Gamma\underline{\unsure{\mathcal G|\widetilde G}}^{\unsure{\natural|\sharp}}\inserted{=\Gamma\underline M},\quad
\underline\mu\in\Gamma\underline{\mathbb O}_{12}^{*}
\]
une section \inserted{locale} de
$\underline G_{\rho,\sigma}^{\unsure{\natural|\sharp}}
=\mathbb O_{12}^{*}\underline{\unsure{\mathcal G|\widetilde G}}^{\unsure{\natural|\sharp}}$,
écrivons qu'elle est dans $\underline N_\rho$ d'où
\margin{78}

\scanpage{80}

\[
\operatorname{int}(f^{-1})\underbrace{u_{\underline\mu}(\rho)}_{\lambda_1^\nu\rho}
=\rho^{\varepsilon_3}\qquad\text{i.e.}\qquad
\lambda_1^\nu\rho=\operatorname{int}(f)\rho^{\varepsilon_3}
\]
en réduisant mod $\underline M$ on trouve
\[
\rho=\deleted{\unsure{\operatorname{int}f}}\rho^{\varepsilon_3}\qquad\text{dans }H^+
\]
ce qui implique $\varepsilon_3=1$, OK. Donc on a
\begin{samepage}
\[
(163)
\]
\diagram{\begin{tikzcd}[ampersand replacement=\&,column sep=large]
\deleted{\underline Z_\rho^{\unsure{\natural|\sharp}},}\underbrace{\underline N_{\rho,\deleted{\unsure{\sigma}}}^{\unsure{\natural|\sharp}}}_{=\underline Z_\rho^{\unsure{\natural|\sharp}}}\arrow[r,"\sim"] \& \underline\gamma'_{\rho,\sigma}=\underline{\mathbb O}_{12}^{*\prime}
\end{tikzcd}}
\end{samepage}
--- ainsi $\inserted{\underline Z_\rho^{\unsure{\natural|\sharp}}=}\underline N_{\rho,\deleted{\unsure{\sigma}}}^{\unsure{\natural|\sharp}}$ est une \uline{section partielle}, au dessus de $\underline{\mathbb O}_{12}^{\unsure{\prime*\prime|*\prime}}$, de l'extension $N_{\rho,\sigma}$ de $\underline{\mathbb O}_{12}^{*}$ par $L_\rho$.

\deleted{\unsure{Posons} donc} Cette section partielle ne
\[
\deleted{(164)\qquad N_{\rho,\sigma}^{?}=\{1,\tau'\}.}
\]
contient pas $\tau'$, qui est une section de $\underline N_{\rho\deleted{\unsure{,\sigma}}}$. Est-ce que $\tau'$ normalise $\underline{\unsure{Z|N}}_\rho^{\unsure{\natural|\sharp}}$ --- comme il normalise $\underline N_\rho$, il suffirait qu'elle normalise
$\deleted{\unsure{\underline{\mathcal G}^{\natural}}}\inserted{\underline G_{\rho,\sigma}^{\unsure{\natural|\sharp}}}$,
\deleted{Mais \inserted{ce qui} est évident car $\underline{\unsure{\mathcal G|\widetilde G}}^{\unsure{\natural|\sharp}}=\underline{\unsure{\mathcal G|\widetilde G}}^0$. Donc $\tau'=\tau\sigma$, où $\tau$}

Comme $\tau'=\tau\sigma$, \deleted{il revient} et
$\tau\in\Gamma\underline{\unsure{G|\mathcal G}}_{\rho,\sigma}^{\unsure{\natural|\sharp}}$
normalise $\underline{\unsure{G|\mathcal G}}_{\rho,\sigma}^{\unsure{\natural|\sharp}}$,
il revient au $\overline{\mathrm m}$ de dire que $\sigma$ normalise
$\underline{\unsure{G|\mathcal G}}_{\rho,\sigma}^{\unsure{\natural|\sharp}}
=\underline{\mathbb O}_{12}^{*}\cdot\underline{\unsure{\mathcal G|\widetilde G}}^{\unsure{\natural|\sharp}}_{\deleted{\unsure{\rho,\sigma}}}$.
Comme $\sigma$ normalise $\underline{\unsure{\mathcal G|\widetilde G}}^{\unsure{\natural|\sharp}}=\underline{\unsure{\mathcal G|\widetilde G}}^0$ qui est distingué, il suffit de regarder l'action de $\sigma$ sur $\underline{\mathbb O}_{12}^{*}\subset G_{\rho,\sigma}$, \deleted{\unsure{défini|définis}} i.e. sur les $u_{\underline\mu}$ de (151). La question est si les
\[
\sigma u_{\underline\mu}\sigma^{-1}u_{\underline\mu}^{-1}\in\Gamma\underline{\unsure{\mathcal G|\widetilde G}}^{\unsure{\natural|\sharp}},
\]
or
\margin{79}

\scanpage{81}
\mspage{688}

\[
\sigma u_{\underline\mu}\sigma^{-1}u_{\underline\mu}^{-1}
=\sigma u_{\underline\mu}(\sigma)^{-1}
=\sigma(\omega^{\nu_2}\sigma)^{-1}=\omega^{\nu_2}
\]
donc

\uline{Proposition}. Pour une section \inserted{$x$} de $\underline G_{\rho,\sigma}^{\unsure{\natural|\sharp}}$, \deleted{$\unsure{\sigma}$}
$\sigma(x)$ est une section de $\underline G_{\rho,\sigma}^{\unsure{\natural|\sharp}}\cdot\underbrace{\{1,\omega\}_{S_0}}_{\mu_{S_0}}$ et son image dans $\Gamma\{1,\omega\}$ est égale à $\nu_2\deleted{\unsure{\chi}}(x)$ --- donc elle est dans $\underline G_{\rho,\sigma}^{\unsure{\natural|\sharp}}$ ssi $\underline\mu(x)$ est une section de
\[
(164)\qquad
\underline{\mathbb O}_{12}^{*\prime\prime}
\overset{\unsure{\mathrm{def}|\mathrm{dfn}}}{=}
\{\underline\mu\in\Gamma\mathbb O_{12}^*\mid\varepsilon_2(\underline\mu)=1\}
\]

\uline{Corollaire 1}. $\underline G_{\rho,\sigma}^{\unsure{\natural|\sharp}}$ n'est pas \unsure{normalisé} par $\sigma$ ni $\tau'$, mais
\[
(165)\qquad
\underline G_{\rho,\sigma}^{\prime\prime\unsure{\natural|\sharp}}
\overset{\mathrm{dfn}}{=}
\{x\in\Gamma\underline G_{\rho,\sigma}^{\unsure{\natural|\sharp}}\mid\varepsilon_2(x)=1\}
\]
est égal à
\[
\underline G_{\rho,\sigma}^{\unsure{\natural|\sharp}}
\cap\sigma(\underline G_{\rho,\sigma}^{\unsure{\natural|\sharp}})
=\underline G_{\rho,\sigma}^{\unsure{\natural|\sharp}}
\cap\tau'(\underline G_{\rho,\sigma}^{\unsure{\natural|\sharp}})
\]
et est normalisé par $\sigma$ et $\tau'$.

\uline{Corollaire 2}. $\underline{\unsure{Z|N}}_\rho^{\unsure{\natural|\sharp}}$ n'est pas normalisé par $\sigma$ ni $\tau'$, mais
\[
(166)\qquad
\underline{\unsure{Z|N}}_\rho^{0\unsure{\natural|\sharp}}
\overset{\mathrm{dfn}}{=}
\underline{\unsure{Z|N}}_\rho^{\unsure{\natural|\sharp}}
\cap\underline G_{\rho,\sigma}^{\prime\prime\unsure{\natural|\sharp}}
\deleted{\unsure{=|\subset}}
\subset\underline{\unsure{Z|N}}_\rho^{\unsure{\natural|\sharp}}
\deleted{\unsure{\underline N_\rho}}
\]
\margin{\uline{NB} On a $N_\rho^{\prime\prime\unsure{\natural|\sharp}}=(N_\rho^{\unsure{\natural|\sharp}})^0$ (comp. neutre)}
est normalisé par $\sigma$ et $\tau'$ --- \deleted{\unsure{s'identifie} donc}
\begin{samepage}
\[
(167)
\]
\diagram{\begin{tikzcd}[ampersand replacement=\&,column sep=large,row sep=small]
\underline{\unsure{Z|N}}_\rho^{0\unsure{\natural|\sharp}}\arrow[r,"\chi_\rho","\sim"']\arrow[d,equal,sloped] \& \underline{\mathbb O}_{12}^{*0}\quad(=\underline{\mathbb O}_{12}^{*\prime}\cap\underline{\mathbb O}_{12}^{*\prime\prime}) \\
(\underline{\unsure{Z|N}}_\rho^{\unsure{\natural|\sharp}})^0 \&
\end{tikzcd}}
\end{samepage}
\margin{A-t-on $\underline Z_\rho^{\prime\prime\unsure{\natural|\sharp}}=(\underline N_\rho)^0$ (comp. neutre) ? Il faudrait aussi \deleted{\unsure{relations}} $\underline G_{\rho,\sigma}^{0}$, $\underline G_{\rho,\sigma}/\underline G_{\rho,\sigma}^{0}$\ldots\ comparer \unsure{avec} les diagrammes (38) p. 622'.}

\uline{Corollaire 3}. Soit
\[
(168)\qquad
\underline N_{\rho(-\bar j)}^{\prime\prime\prime\deleted{\unsure{\natural|\sharp}}}
=\underbrace{\{1,\tau'\}_{S_0}}_{(\mu_{\tau'})_{S_0}}
\underbrace{\underline{\unsure{Z|N}}_\rho^{0\unsure{\natural|\sharp}}}_{(\underline{\unsure{Z|N}}_\rho^{\unsure{\natural|\sharp}})^0}
\subset\underline N_\rho^{\deleted{\unsure{\natural|\sharp}}},
\]
alors on a
\begin{samepage}
\[
(169)
\]
\diagram{\begin{tikzcd}[ampersand replacement=\&,column sep=large]
\underline N_{\rho(-\bar j)}^{\prime\prime\prime\deleted{\unsure{\natural|\sharp}}}
\arrow[r,"\chi_\rho","\sim"'] \&
\underline{\mathbb O}_{12}^{*\prime\prime\prime}
=\{\underline\mu\in\Gamma\mathbb O_{12}^*\mid(\varepsilon_2\varepsilon_3)(\underline\mu)=1\}
\end{tikzcd}}
\end{samepage}
\margin{80}

\scanpage{82}

Donc \(N_\rho^{'''\unsure{\natural|\sharp}}\) est une section partielle de \(N_\rho\)
sur \(\mathbb O_{12}^{*'''}\). Enfin, posons
\[
(170)\qquad
\underline Z_\rho^{!}
=\underline Z_\rho^{0\unsure{\natural|\sharp}}\times\mu_{S_0}
\qquad(\mu=\{1,\omega\})
\]
alors \(\underline Z_\rho^!\) est normalisé par \(\tau'\), et
on peut introduire
\[
(171)\qquad
\deleted{\underline{\unsure{Z}}}\inserted{\underline N}{}_\rho^{?}
\mathrel{\overset{\mathrm{dfn}}{=}}
(\unsure{\mu_{\tau'}|\mu\tau'})_{S_0}\cdot\underline Z_\rho^!
\subset\underline N_\rho
\]
de sorte que \(\underline Z_\rho^{0\unsure{\natural|\sharp}}\) est d'indice \(4\)
dans \(\underline N_\rho^{?}\). On a alors

\uline{Corollaire 4.} On a une suite exacte
\diagram{
\begin{tikzcd}[ampersand replacement=\&,column sep=large]
{(172)\quad 1}
  \arrow[r] \&
{\underset{\substack{\parallel\,\mathrm{dfn}\\ \mu_\rho}}{\mu_{S_0}}}
  \arrow[r] \&
{\underline N_\rho^{?}}
  \arrow[r] \&
{\underbrace{\underline{\mathbb O}_{12}^{*}}_{\gamma_{\rho,\sigma}}}
  \arrow[r] \& 1
\end{tikzcd}
}
\margin{\uline{NB} Il n'y a pas lieu de \unsure{redésigner|distinguer}
\(\mathbb O_{12}^{*}\) --- \unsure{il n'y a pas},
\(\unsure{N_\rho|N_\rho^{!}}\) \unsure{etc}, \unsure{aucun point de}
\unsure{risque de confusion} \unsure{des groupes} \illegible{1}}
\(\underline N_\rho^{?}\) est une bi-section de l'extension
\(\underline N_\rho\) de \(\underline{\mathbb O}_{12}^{*}\) par \(L_\rho\).

Il y aurait lieu de développer le diagramme des sous-groupes
remarquables de \(\underline N_\rho\), calqué sur
le diag. \((38)\) des p.\(622^{\unsure{1|'}}\).

\margin{81}

\scanpage{83}
\mspage{689}

8°) \uline{Les groupes} \(N_\sigma,N_\sigma^{\unsure{\natural|\sharp}}\ldots\).

Posons
\[
(173)\qquad
\left\{
\begin{gathered}
\underline N_\sigma\subset\underline{\unsure{G}}_{\rho,\sigma}\\
\mathrel{\overset{\mathrm{dfn}}{=}}
\left\{\underline u\in\Gamma\underline G_{\rho,\sigma}\ \middle|\
\underline u(\sigma)=
\deleted{\varepsilon_2(\underline u)}
\sigma^{\varepsilon_2(\underline u)}
\left(=\omega^{\nu_2(\underline u)}\sigma\right)\right\}.
\end{gathered}
\right.
\]
\[
(174)\qquad
\underline N_\sigma^{\unsure{\natural|\sharp}}
\mathrel{\overset{\mathrm{dfn}}{=}}
\underline N_\sigma\cap\underline G_{\rho,\sigma}^{\unsure{\natural|\sharp}}
\]
L'hom composé
\diagram{
\begin{tikzcd}[ampersand replacement=\&,column sep=large]
\underline N_\sigma^{\unsure{\natural|\sharp}}
  \arrow[r,hook,"{\deleted{\unsure{\underline\chi_\sigma}}}"]
  \arrow[rr,bend right=20,"{\underline\chi_\sigma}"'] \&
{\deleted{\underline{\unsure{\mathbb O}}}\inserted{\underline G}{}_{\rho,\sigma}^{\unsure{\natural|*}}}
  \arrow[r,"{\underline\chi}"] \&
\mathbb O_{12}^{*}
\end{tikzcd}
}
est épi. En effet, les \(u_{\underline\mu}\) sont dans
\deleted{\(\unsure{G}\)}
\(N_{\unsure{\sigma|\rho}}^{\unsure{\natural|\sharp}}\) --- donc on a $\overline{\mathrm m}$ un scindage
\diagram{
\begin{tikzcd}[ampersand replacement=\&,column sep=large]
\mathbb O_{12}^{*}\arrow[r,hook] \&
\underline N_\sigma^{\unsure{\natural|\sharp}}
\end{tikzcd}
},
d'où des décompositions en produits \(\frac12\) directs
\[
(175)\qquad
\left\{
\begin{aligned}
\underline N_\sigma^{\unsure{\natural|\sharp}}
&\simeq\mathbb O_{12}^{*}\cdot S\underline Z_\sigma^{\unsure{\natural|\sharp}}\\
\underline N_\sigma
&\simeq\mathbb O_{12}^{*}\cdot S\underline Z_\sigma
\end{aligned}
\right.
\]
où on a posé
\[
(176)\qquad
\left\{
\begin{aligned}
\underline Z_\sigma
&=\operatorname{Ker}\left(
\underline N_\sigma\xrightarrow{\varepsilon_2\circ\underline\chi_\sigma}\mu_{S_0}
\right)
=\underline{\operatorname{Cent}}_{\underline G_{\rho,\sigma}}(\sigma)\\
S\underline Z_\sigma
&=\underline Z_\sigma\cap\underline{\unsure{\mathcal G|\widetilde G}}^{+}
=\underline{\operatorname{Cent}}_{\underline{\unsure{\mathcal G|\widetilde G}}^{+}}(\sigma)\\
\underline Z_\sigma^{\unsure{\natural|\sharp}}
&=\underline Z_\sigma\cap\underline G_{\rho,\sigma}^{\unsure{\natural|\sharp}}\\
S\underline Z_\sigma^{\unsure{\natural|\sharp}}
&=S\underline Z_\sigma\cap
\underbrace{
\deleted{\underline G_{\rho,\sigma}}\inserted{\underline{\unsure{\mathcal G|\widetilde G}}}{}^{\unsure{\natural|\sharp}}
}_{=\underline{\unsure{\mathcal G|\widetilde G}}^0=\underline M}
\deleted{
=\underline{\unsure{Z}}_\sigma\cap
\underbrace{\underline{\unsure{\mathcal G|\widetilde G}}^{\unsure{\natural|\sharp}}}
_{\parallel\,\underline{\unsure{\mathcal G|\widetilde G}}^0=\underline M}
}
\end{aligned}
\right.
\]
D'ailleurs on a
\diagram{
\begin{tikzcd}[ampersand replacement=\&,column sep=large,row sep=large]
S\underline Z_\sigma\arrow[r] \&
{\deleted{\unsure{H}_{S_0}}\underline{\operatorname{Cent}}_{H^+_{S_0}}(\sigma)}\\
\&
{\underline L_\sigma\simeq(\mathbb Z/4\mathbb Z)_{S_0}}
\arrow[u,"{\unsure{2|\wr}}","{\unsure{\delta|\wr}}"']
\arrow[ul,hook,dashed,"{\unsure{\hookleftarrow|\leftrightarrow}}"']
\end{tikzcd}
}
qui a une section, donc
\[
(177)\qquad
S\underline Z_\sigma\simeq L_\sigma\times S\underline Z_\sigma^{\unsure{\natural|\sharp}}
\]

\margin{82}

\scanpage{84}

Enfin
\[
(178)\qquad
\begin{gathered}
S\underline Z_\sigma^{\unsure{\natural|\sharp}}\subset\underline M\\
\parallel\\
\left\{x\in\Gamma\underline M\ \middle|\ (1-\sigma_{\underline M})(x)=0\right\}
=\left\{
\deleted{\unsure{x=a_0\lambda_0+a_1\lambda_1}}
\inserted{a_0\lambda_0+a_1\lambda_1}
\ \middle|\ a_0=a_1
\right\}\\
=\left\{a(\lambda_0+\lambda_1)\ \middle|\ a\in\Gamma\mathbb G_a\right\}
\simeq\mathbb G_a,
\end{gathered}
\]
d'où, compte tenu de \(177\) :
\[
(179)\qquad
\left\{
\begin{aligned}
S\underline Z_\sigma
&\simeq\underline L_\sigma\times
\underbrace{\mathbb G_a}_{S\underline Z_\sigma^{\unsure{\natural|\sharp}}}\\
S\underline Z_\sigma^{\unsure{\natural|\sharp}}
&=(S\underline Z_\sigma)^0
\quad\text{(comp. neutre)}\simeq\mathbb G_a
\end{aligned}
\right.
\]
On a donc
\diagram{
\begin{tikzcd}[ampersand replacement=\&,column sep=large]
{\deleted{(180)}\quad 1}
 \arrow[r] \&
{\underset{\substack{\rotatebox{90}{$\simeq$}\\S\underline Z_\sigma^{\unsure{\natural|\sharp}}}}{\mathbb G_a}}
 \arrow[r] \&
\underline{\unsure{N}}_\sigma^{\unsure{\natural|\sharp}}
 \arrow[r] \&
{\underset{\substack{\parallel\\\gamma_{\rho,\sigma}}}{\mathbb O_{12}^{*}}}
 \arrow[r] \& 1
\end{tikzcd}
},
plus précisément les isom. \((175)(177)(179)\) donnent
\[
(180)\qquad
\begin{aligned}
\underline{\unsure{N}}_\sigma^{\unsure{\natural|\sharp}}
&\simeq\mathbb O_{12}^{*}\cdot\mathbb G_a\\
\underline N_\sigma
&\simeq\mathbb O_{12}^{*}\cdot
(\underline L_\sigma\times\mathbb G_a)
\end{aligned}
\]
Dans ces formules, \(\mathbb O_{12}^{*}\) opère sur \(\mathbb G_a\)
via
\diagram{
\begin{tikzcd}[ampersand replacement=\&]
{\chi:\mathbb O_{12}^{*}}\arrow[r,dash] \& \mathbb G_m
\end{tikzcd}
},
et sur \(L_\sigma\) via
\diagram{
\begin{tikzcd}[ampersand replacement=\&]
{\varepsilon_2:\mathbb O_{12}^{*}}\arrow[r] \&
{(\pm1)_{S_0}}
\end{tikzcd}
} :
\[
(181)\qquad
\begin{gathered}
\deleted{\underline\mu(\sigma)=\sigma^{\varepsilon_2(\underline\mu)}}\\
\left\{
\begin{aligned}
\underline\mu(\lambda)&=\mu\cdot\lambda\\
\underline\mu(\sigma)&=\sigma^{\varepsilon_2}
\end{aligned}
\right.
\qquad\text{où }\mu=\chi(\underline\mu),\ \varepsilon_2=\varepsilon_2(\underline\mu)
\end{gathered}
\]

\margin{83}

\scanpage{85}
\mspage{690}

9°) \uline{Retour sur} \(\underline M_{\rho,\sigma}\).

Nous inspirant des développements du
\(\S\,49\) (notamment \uline{VI}, 5°)), nous \unsure{poserons}
\[
(182)\qquad
\left\{
\begin{aligned}
\unsure{S_0|s_0}\underline M_{\rho,\sigma}
&=\underline Z_{\rho,\sigma}
\times_\gamma\underline N_\rho^{?}
\times_\gamma\underline N_\sigma^{\unsure{\natural|\sharp}}
\times_\gamma\underline G_{\rho,\sigma}\\
\unsure{S_0|s_0}\underline M_{\rho,\sigma}^{\unsure{\natural|\sharp}}
&=\underline Z_{\rho,\sigma}
\times_\gamma\underline N_\rho^{?}
\times_\gamma\underline N_\sigma^{\unsure{\natural|\sharp}}
\times_\gamma\underline G_{\rho,\sigma}^{\unsure{\natural|\sharp}}
\end{aligned}
\right.
\]
\[
(183)\qquad
\left\{
\begin{aligned}
\deleted{\unsure{S_0}}\underline M_{\rho,\sigma}
&=\underline N_\rho^{?}
\times_\gamma\underline N_\sigma^{\unsure{\natural|\sharp}}
\times_\gamma\underline G_{\rho,\sigma}\\
\underline M_{\rho,\sigma}^{\unsure{\natural|\sharp}}
&=\underline N_\rho^{?}
\times_\gamma\underline N_\sigma^{\unsure{\natural|\sharp}}
\times_\gamma\underline G_{\rho,\sigma}^{\unsure{\natural|\sharp}}
\end{aligned}
\right.
\]
\[
(184)\qquad
\left\{
\begin{aligned}
\unsure{S_0|s_0}\Gamma_{\rho,\sigma}
&=\underline Z_{\rho,\sigma}
\times_\gamma\underline N_{\unsure{\rho|1}}^{?}
\times_\gamma\underline N_\sigma^{\unsure{\natural|\sharp}}\\
\Gamma_{\rho,\sigma}
&=\underline N_\rho^{?}
\times_\gamma\underline N_\sigma^{\unsure{\natural|\sharp}}
\end{aligned}
\right.
\]
où \deleted{\(\unsure{\gamma=}\)}
\[
\gamma=\gamma_{\rho,\sigma}=\mathbb O_{12}^{*}.
\]

En fait, les groupes
\(\unsure{S_0|s_0}\underline M_{\rho,\sigma},
\unsure{S_0|s_0}\underline M_{\rho,\sigma}^{\unsure{\natural|\sharp}}\),
\(\unsure{S_0|s_0}\Gamma_{\rho,\sigma}\)
jouent plutôt un rôle technique auxiliaire, ce sont des extensions des
groupes \(\underline M_{\rho,\sigma}\),
\(\underline M_{\rho,\sigma}^{\unsure{\natural|\sharp}}\)
ou \(\underline\Gamma_{\rho,\sigma}\) par
\[
L_0^{\underline M}
=\operatorname{Ker}\left(
\underline Z_{\rho,\sigma}\longrightarrow
\underset{\substack{\parallel\\\gamma}}{\mathbb O_{12}^{*}}
\right),
\]
\inserted{les deux premières}
déduites
de l'extension
\(\unsure{S_0|s_0}\underline\Gamma_{\rho,\sigma}\)
de \(\underline\Gamma_{\rho,\sigma}\) par
\(L_0^{\underline M}\simeq\mathbb G_a\),
par image inverse via les épis
\(\underline M_{\rho,\sigma}\longrightarrow\underline\Gamma_{\rho,\sigma}\)
et l'épi induit
\(\underline M_{\rho,\sigma}^{\unsure{\natural|\sharp}}
\longrightarrow\underline\Gamma_{\rho,\sigma}\).

Nous \unsure{attacherons} notre attention surtout
aux autres groupes (sans \(\unsure{S_0|s_0}\)), reliés
par un diagramme de suites
exactes (\unsure{comparer} p.\(639\) \((79)\)) :

\margin{84}

\scanpage{86}

\diagram{
\begin{tikzcd}[ampersand replacement=\&,column sep=large,row sep=small,nodes in empty cells]
\deleted{\unsure{1}}
\arrow[r,phantom,"{\deleted{\unsure{\longrightarrow}}}" description]
\& \deleted{\underline{\unsure{M}}^{\unsure{\natural|\sharp}}}
\arrow[r,phantom,"{\deleted{\longrightarrow}}" description] \& {}\\
\&\deleted{\underline{\unsure{\mathcal G|\widetilde G}}} \& {}\\
\&\deleted{\underline{\unsure{\mathcal G|\widetilde G}}^{\unsure{\natural|\sharp}}} \& {}
\end{tikzcd}
}

\diagram{
\begin{tikzcd}[ampersand replacement=\&,column sep=large,row sep=large]
{}\arrow[d,phantom,"{(190)\quad\left\{\vphantom{\begin{matrix}\mathcal G\\[3ex]\mathcal G\end{matrix}}\right.}" description]
\&1\arrow[r,dash]
\&\underbrace{\underline{\unsure{\mathcal G|\widetilde G}}^{\unsure{\natural|\sharp}}}_{\underline M}\arrow[r]\arrow[d,hook]
\&\underline{\unsure{M}}_{\rho,\sigma}^{\unsure{\natural|\sharp}}\arrow[r]\arrow[d,hook]
\&\underline\Gamma_{\rho,\sigma}\arrow[r]\arrow[d,"1"]\&1\\
{}\&1\arrow[r]
\&\underline{\unsure{\mathcal G|\widetilde G}}^{+}\arrow[r]
\&\underline{\unsure{M}}_{\rho,\sigma}\arrow[r,dash]
\&\underline\Gamma_{\rho,\sigma}\arrow[r,dash]\&1
\end{tikzcd}
}

et \deleted{des suites exactes}

\diagram{
\begin{tikzcd}[ampersand replacement=\&,column sep=large,row sep=large]
(191)\&1\arrow[r]
\&S\underline N_\rho^{?}\times
S\deleted{\underline{\unsure{N}}}\inserted{\underline Z}_{\deleted{\unsure{\rho}}\sigma}^{\unsure{\natural|\sharp}}\times
\underset{\underline M}{\underline{\unsure{\mathcal G|\widetilde G}}^{\unsure{\natural|\sharp}}}
\arrow[r]\arrow[d,hook]
\&\underline M_{\rho,\sigma}^{\unsure{\natural|\sharp}}\arrow[r]\arrow[d,hook]
\&\underset{\deleted{\unsure{\gamma}}}{\mathbb O_{12}^{*}}
\arrow[r]\arrow[d,phantom,"{\unsure{\Downarrow|\Vert}}" description]
\&1\\
{}\&1\arrow[r]
\&S\underline N_\rho^{?}\times
S\deleted{\underline{\unsure{N}}}\inserted{\underline Z}_\sigma^{\unsure{\natural|\sharp}}\times
\underline{\unsure{\mathcal G|\widetilde G}}^{+}\arrow[r]
\&\underline M_{\rho,\sigma}\arrow[r,dash]
\&\underbrace{\mathbb O_{12}^{\unsure{*|\wedge}}}_{\gamma_{\rho,\sigma}}\arrow[r,dash]
\&1
\end{tikzcd}
},

\diagram{
\begin{tikzcd}[ampersand replacement=\&,column sep=large]
(192)\&1\arrow[r,dash]
\&S\underline N_\rho^{?}\times
S\deleted{\underline{\unsure{N}}}\inserted{\underline Z}_\sigma^{\unsure{\natural|\sharp}}\arrow[r]
\&\underline\Gamma_{\rho,\sigma}\arrow[r,dash]
\&\underbrace{\mathbb O_{12}^{\unsure{*|\wedge}}}_{\gamma_{\rho,\sigma}}\arrow[r]
\&1
\end{tikzcd}
},

où on a \deleted{posé (}
\[
(193)\qquad
\left\{
\begin{aligned}
S\underline N_\rho^{?}
&=\mu_\rho\simeq\{\pm1\}_{S_0}&& (172)\\
S Z_\sigma^{\unsure{\natural|\sharp}}
&\simeq\mathbb G_a&&(179)
\end{aligned}
\right.
\quad .
\]

Ces suites exactes se précisent, grâce
aux décompositions \(1/2\) directes :
\[
\deleted{
(194)\qquad
\left\{
\begin{aligned}
\underline{\unsure{Z|G}}_{\rho,\sigma}
&\simeq\mathbb O_{12}^{\unsure{*|\wedge}}\cdot\mathbb G_a
&&(150,151)\\
\unsure{S}\underline{\unsure{N|Z}}_\sigma^{\unsure{\natural|\sharp}}
&\simeq\mathbb O_{12}^{*}\cdot\mathbb G_a
&&(180)\\
N_\rho^{?}/\mu_\rho
&\unsure{\simeq}\simeq\mathbb O_{12}^{*}
&&(172)
\end{aligned}
\right.
},
\]
\[
(194)\qquad
\left\{
\begin{aligned}
\underline Z_{\rho,\sigma}
&\simeq\mathbb O_{12}^{\unsure{*|\wedge}}\cdot\mathbb G_a
&&(150,151)\\
\underline G_{\rho,\sigma}
&\simeq\mathbb O_{12}^{*}\cdot
\underline{\unsure{\mathcal G|\widetilde G}}^{+}
&&\text{id}\\
\underline G_{\rho,\sigma}^{\unsure{\natural|\sharp}}
&\simeq\mathbb O_{12}^{\unsure{*|\wedge}}\cdot
\underbrace{
\deleted{\underline{\unsure{\mathcal G|\widetilde G}}}
\inserted{\underline M}\deleted{{}^{\unsure{\natural|\sharp}}}
}_{\underline{\unsure{\mathcal G|\widetilde G}}^{\unsure{\natural|\sharp}}}
&&(160)\\
\underline N_\sigma^{\unsure{\natural|\sharp}}
&\simeq\mathbb O_{12}^{*}\cdot\mathbb G_a
&&(180)\\
&\text{et enfin}\\
\underline N_\rho/\mu_\rho
&\simeq\mathbb O_{12}^{*}
&&(172)
\end{aligned}
\right.
\]

\margin{85}

\scanpage{87}
\mspage{691}

On trouve alors des isom. canoniques

\[
(195)\quad
\left\{\vcenter{\hbox{\diagram{
\begin{tikzcd}[ampersand replacement=\&,column sep=small,row sep=large]
\underline M_{\rho,\sigma}/\unsure{\mu_\rho}
 \arrow[r,phantom,"\simeq"]
 \& \mathbb O_{12}^{*}\cdot
   (\mathbb G_a\times\underline{\unsure{\mathcal G|\widetilde G}}^{+})
\\
\underline M_{\rho,\sigma}^{\unsure{\natural|\sharp}}/\mu_\rho
 \arrow[u,hook] \arrow[r,phantom,"\simeq"]
 \& \mathbb O_{12}^{*}\cdot
   \bigl(\underbrace{\mathbb G_a\times\underline M}_{
   \unsure{\mathbb G_a\times\underline{\unsure{\mathcal G|\widetilde G}}^{+}}}\bigr)
   \arrow[u,equal,xshift=-2em] \arrow[u,hook,xshift=4em]
\\
\underline\Gamma_{\rho,\sigma}\arrow[r,phantom,"\simeq"]
 \& \mathbb O_{12}^{*}\cdot\mathbb G_a
\end{tikzcd}
}}}\right.
\]

\margin{op. de $\mathbb O_{12}^{*}$ sur
$\mathbb G_a\times\underline{\unsure{\mathcal G|\widetilde G}}^{+}$ \unsure{données}
par (181) et (151). Sur $\mathbb G_a\times\underline M$
\deleted{\unsure{elle}}, $\mathbb O_{12}^{*}$
opère simplement par \unsure{homothétie}.}

$\underline M_{\rho,\sigma}$ et
$\underline M_{\rho,\sigma}^{\unsure{\natural|\sharp}}$,
\unsure{en tant qu'extensions}, de
$\underline M_{\rho,\sigma}/\mu_\rho$ et
$\underline M_{\rho,\sigma}^{\unsure{\natural|\sharp}}/\mu_\rho$
par $\mu_\rho$, se déduisent de l'extension

\[
(196)\quad\vcenter{\hbox{\diagram{
\begin{tikzcd}[ampersand replacement=\&,column sep=small]
1\arrow[r]\&\mu_\rho\arrow[r]\&N_\rho^{?}\arrow[r]
\&\mathbb O_{12}^{*}\arrow[r]\&1
\end{tikzcd}
}}}
\]

de $\mathbb O_{12}^{*}$ par $\mu_\rho$,
par images inverses des hom. canoniques des groupes (195)
sur le facteur $\mathbb O_{12}^{*}$.
Il faudrait revenir sur l'extension (196), qui est sûrement
l'image inverse d'une \unsure{ext} de
$\unsure{(\mu\times\mu)_{S_0}}$ par $\mu$, via

\diagram{
\begin{tikzcd}[ampersand replacement=\&]
\underline{\mathbb O}_{12}^{*}
 \arrow[r,"{(\varepsilon_2,\varepsilon_3)}"]
 \&(\mu\times\mu)_{S_0}
\end{tikzcd}
}

(Comparer \S\,49 \uline{$\mathrm V$} $8^{\circ}$).
\deleted{L'opération}

L'opération de $\mathbb O_{12}^{*}$ sur
$\mathbb G_a\times\underline{\unsure{\mathcal G|\widetilde G}}^{+}$
(celle sur $\mathbb G_a\times\underline{\unsure{\mathcal G|\widetilde G}}^{\unsure{\natural|\sharp}}
=\mathbb G_a\times\underline M$) se fait via des opérations sur
$\mathbb G_a$ et $\underline{\unsure{\mathcal G|\widetilde G}}^{+}$, qui sont explicitées
\deleted{dans} par (181) et (151).

\margin{86}

\scanpage{88}

$10^{\circ})$ \uline{Les sous-groupes}
$\underline M_{\rho,\sigma}[0]$, $\underline M_{\rho,\sigma}(0)$,
$\underline M_{\rho,\sigma}(\frac12)$,
$\underline M'_{\rho,\sigma}(-\bar\jmath)$,
$[M'''_{\rho,\sigma}(-\bar\jmath),M^{?}_{\rho,\sigma}(-\bar\jmath)]$.

On utilise le \unsure{plongement|pl-gement}

\[
(195)\quad
\left\{\begin{array}{l}
\underline M_{\rho,\sigma}
=\underline N_\rho^{?}\times_\gamma
 \underline N_\sigma^{\unsure{\natural|\sharp}}\times_\gamma
 \underline G_{\rho,\sigma}
\subset\underline N_\rho^{?}\times
 \underline N_\sigma^{\unsure{\natural|\sharp}}\times
 \underline G_{\rho,\sigma}
\\[-2pt]
\quad\Vert
\\[-2pt]
\bigl\{(a,b,U)\,\unsure{\{}\in
 \Gamma\deleted{\unsure{=}}\underline N_\rho^{?}
 \times\underline N_\sigma\times\underline G_{\rho,\sigma}
 \mid\underline\chi_\rho(a)\unsure{\equiv|=}
 \underline\chi_\sigma(b)\unsure{\equiv|=}
 \deleted{\unsure{\chi}}\underline\chi_G(U)\bigr\}
\end{array}\right.
\]

On définit donc

\[
(196)\quad\begin{gathered}
\underline M_{\rho,\sigma}[0]\subset
\underline M_{\rho,\sigma}^{\deleted{\unsure{\natural|\sharp}}}
\\[-2pt]
\Vert
\\[-2pt]
\bigl\{(a,b,U)\in\Gamma\underline M_{\rho,\sigma}
 \mid U\in Z_{\rho,\sigma}\bigr\}
\end{gathered}
\]

\margin{NB on aura $\overline{\mathrm m}$
$\underline M_{\rho,\sigma}[0]\subset
\underline M_{\rho,\sigma}^{\unsure{\natural|\sharp}}$.}

donc
\diagram{
\begin{tikzcd}[ampersand replacement=\&,column sep=small]
\underline M_{\rho,\sigma}[0]
\simeq\underline N_\rho^{?}\times_\gamma
\underline N_\sigma^{\unsure{\natural|\sharp}}\times_\gamma Z_{\rho,\sigma}
 \arrow[r,"\sim"]
 \&\deleted{\unsure{\underline Z}}\unsure{S_0|s_0}\underline\Gamma_{\rho,\sigma}
\end{tikzcd}
}

\deleted{Posons}\inserted{Soit}

\[
(197)\quad\underline Z_{\rho,\sigma}(0)\subset\underline Z_{\rho,\sigma}
\]

le facteur $\frac12$ direct $\mathbb O_{12}^{*}$
de (150) (151),

\margin{\uline{NB}
$\underline Z_{\rho,\sigma}(0)\subset\underline Z_{\rho,\sigma}
\subset\underline G_{\rho,\sigma}$
\unsure{est aussi} le fact. $\frac12$ direct
$\mathbb O_{12}^{*}$ de $\underline G_{\rho,\sigma}$
dans (150), il pourrait \unsure{se noter} aussi
$\underline G_{\rho,\sigma}(0)$.}

\deleted{Donc}\inserted{On a} alors, posons
\inserted{$\underline M_{\rho,\sigma}(0)\simeq
\underline N_\rho^{?}\times_\gamma
\underline N_\sigma^{\unsure{\natural|\sharp}}\times_\gamma
Z_{\rho,\sigma}(0)$}

\[
(198)\quad
\left\{\vcenter{\hbox{\diagram{
\begin{tikzcd}[ampersand replacement=\&,column sep=large,row sep=large]
\underline M_{\rho,\sigma}(0)
 \arrow[r,phantom,"\subset"]\arrow[d,"\sim"]
 \&\underline M_{\rho,\sigma}[0]\arrow[d,"\sim"]
\\
\unsure{S\underline\Gamma_{\rho,\sigma}(0)}
 \arrow[r,hook]\arrow[d,"\sim"]
 \&\unsure{S_0|s_0}\underline\Gamma_{\rho,\sigma}
\\
\underline\Gamma_{\rho,\sigma}(0)\&
\end{tikzcd}
}}}\right.
\]

\[
(199)\quad\begin{aligned}
\underline M_{\rho,\sigma}(0)
&=\bigl\{(a,b,U)\in\Gamma\underline M_{\rho,\sigma}
 \mid U\in\Gamma\underline Z_{\rho,\sigma}(0)\bigr\}
\\
&=\underline\mu\cdot
\left(\underset{\substack{\rotatebox{90}{$\in$}\\\Gamma\mathbb G_a}}{\lambda},
\underset{\substack{\rotatebox{90}{$\in$}\\
\Gamma\underline{\unsure{\mathcal G|\widetilde G}}^{+}}}{0}\right)
\end{aligned}
\]

\margin{dans la description (195) de $\underline M_{\rho,\sigma}$
comme prod. $\frac12$ direct.}

\margin{87}

\scanpage{89}
\mspage{692}

\margin{\uline{NB}
\deleted{\unsure{$\underline M_{\rho,\sigma}(0)$}}
$\mu_\rho\subset\underline M_{\rho,\sigma}(0)
\subset\underline M_{\rho,\sigma}[0]$ i.e.}

\[
(200)\quad\vcenter{\hbox{\diagram{
\begin{tikzcd}[ampersand replacement=\&,column sep=small]
\underline M_{\rho,\sigma}
 \deleted{/\mu_\rho}(0)/\mu_\rho
 \simeq\mathbb O_{12}^{*}\cdot\mathbb G_a
 \arrow[r,hook]
 \&\underline M_{\rho,\sigma}\simeq
 \mathbb O_{12}^{*}
 (\mathbb G_a\times\underline{\unsure{\mathcal G|\widetilde G}}^{+})
\end{tikzcd}
}}}
\]

(\inserted{\unsure{correspondant à}}l'\unsure{inclusion} canonique
\deleted{\unsure{)}} de $\mathbb G_a$ dans
$\mathbb G_a\times\underline{\unsure{\mathcal G|\widetilde G}}^{+}$).

D'autre part on aura

\[
(201)\quad\vcenter{\hbox{\diagram{
\begin{tikzcd}[ampersand replacement=\&,column sep=small]
\underline M_{\rho,\sigma}
 \deleted{/\mu_\rho}[0]/\mu_\rho
 \simeq\mathbb O_{12}^{*}\cdot
 (\mathbb G_a\times L_0^{\underline M})\arrow[r,hook]
 \&\overbrace{\mathbb O_{12}^{*}
 (\mathbb G_a\times\underline{\unsure{\mathcal G|\widetilde G}}^{+})}
 ^{M_{\rho,\sigma}\simeq}
\end{tikzcd}
}}}
\]

On définit d'autre part

\[
(202)\quad
\underline M_{\rho,\sigma}(\tfrac12)
=\bigl\{(a,b,U)\in\Gamma\underline M_{\rho,\sigma}\mid U=b\bigr\}
\]

(donc on aura $\underline M_{\rho,\sigma}(\frac12)
\subset\underline M_{\rho,\sigma}^{\unsure{\natural|\sharp}}$)

Donc \unsure{on a|le}

\[
(203)\quad\left\{\begin{array}{l}
\vcenter{\hbox{\diagram{
\begin{tikzcd}[ampersand replacement=\&]
\underline M_{\rho,\sigma}(\frac12)\arrow[r,"\sim"]
 \&\underline\Gamma_{\rho,\sigma}
\end{tikzcd}
}}}
\\[6pt]
\underline M_{\rho,\sigma}^{\unsure{\natural|\sharp}}
 \simeq\underline M_{\rho,\sigma}(\frac12)\cdot\underline M
\\[6pt]
\underline M_{\rho,\sigma}\simeq
 \underline M_{\rho,\sigma}(\frac12)\cdot
 \underline{\unsure{\mathcal G|\widetilde G}}^{+}.
\end{array}\right.
\]

Dans la représentation de $\underline M_{\rho,\sigma}/\mu_\rho$,
$\underline M_{\rho,\sigma}^{\unsure{\natural|\sharp}}/\mu_\rho$
en produits $\frac12$ directs, on aura la description

\margin{\unsure{i.e.} $\mu_\rho\subset\underline M_{\rho,\sigma}(\frac12)$}

\[
(204)\quad\left\{\begin{array}{l}
\underline M_{\rho,\sigma}(\frac12)/\mu_\rho
=\left\{\underline\mu\cdot\left(
\underset{\substack{\rotatebox{90}{$\in$}\\\Gamma\mathbb G_a}}{\lambda},
\underbrace{\lambda(\lambda_1+\lambda_2)}_{
\substack{\rotatebox{90}{$\in$}\\\Gamma\underline M}}\right)
\mathrel{\Big|}\unsure{-\lambda|\lambda}\in\Gamma\mathbb G_a\right\}
\\[-2pt]
\qquad\unsure{\rotatebox{90}{$\simeq$}|\Vert}
\\[-2pt]
\vcenter{\hbox{\diagram{
\begin{tikzcd}[ampersand replacement=\&,column sep=large,row sep=small]
\& |[draw,rounded corners]| M_{\rho,\sigma}^{\unsure{\natural|\sharp}}\simeq
 \arrow[d,no head,bend right=15]
\\
\mathbb O_{12}^{*}\cdot\mathbb G_a
 \arrow[r,hook]
 \&\mathbb O_{12}^{*}(\mathbb G_a\times\underline M)
\end{tikzcd}
}}}
\end{array}\right.
\]

associé à
\diagram{
\begin{tikzcd}[ampersand replacement=\&,row sep=small]
\mathbb G_a\arrow[r]\&\mathbb G_a\times M
\\
\lambda\arrow[r,mapsto]\&\lambda,\lambda(\lambda_1+\lambda_2)
\end{tikzcd}
}

\margin{88}

\scanpage{90}

Enfin on pose

\[
(205)\quad\begin{gathered}
\underline M_{\rho,\sigma}^{\unsure{\deleted{?}|'}}(-\bar\jmath)
\subset\underline M_{\rho,\sigma}
\\[-2pt]
\Vert
\\[-2pt]
\bigl\{(a,b,U)\in\Gamma\underline M_{\rho,\sigma}\mid U=a\bigr\}
\end{gathered}
\]

ici on n'aura pas néc.
$U\in\Gamma(\underline G_{\rho,\sigma}^{\unsure{\natural|\sharp}}
\deleted{\unsure{\underline M}})$ car
\deleted{\illegible{2}}
$\underline N_\rho^{?}\not\subset
\underline G_{\rho,\sigma}^{\unsure{\natural|\sharp}}$, donc

\[
\underline M_{\rho,\sigma}^{\unsure{\deleted{?}|'}}(-\bar\jmath)
\not\subset\underline M_{\rho,\sigma}^{\unsure{\natural|\sharp}},
\]

contrairement à ce qui se passe pour
$\underline M_{\rho,\sigma}(0)$,
$\underline M_{\rho,\sigma}[0]$,
$\underline M_{\rho,\sigma}(\frac12)$.
On aura un

\[
(206)\quad\left\{\begin{array}{l}
\vcenter{\hbox{\diagram{
\begin{tikzcd}[ampersand replacement=\&]
\underline M_{\rho,\sigma}^{\unsure{\deleted{?}|'}}(-\bar\jmath)
 \arrow[r,"\sim"]\&\underline\Gamma_{\rho,\sigma}
\end{tikzcd}
}}}
\\[6pt]
\underline M_{\rho,\sigma}\simeq
\underline M_{\rho,\sigma}(-\bar\jmath)\cdot
\underline{\unsure{\mathcal G|\widetilde G}}^{+}
\end{array}\right.
\]

Mais on \deleted{\unsure{n'aura pas}}\inserted{a}
\[
\underline\mu_\rho\not\subset\underline M_{\rho,\sigma}(-\bar\jmath)
\]

plus précisément
\[
\underline\mu_\rho\cap\underline M_{\rho,\sigma}(-\bar\jmath)=\{1\}
\]

donc \deleted{la relation}
$\underline M_{\rho,\sigma}(-\bar\jmath)$ ne définit pas
(comme $\underline M_{\rho,\sigma}(\frac12)$) une section de
$\underline M_{\rho,\sigma}/\underline\mu_\rho$
\deleted{\unsure{sur}}\inserted{\unsure{au}} dessus de
$\deleted{\unsure{\mu_2}}\underline\Gamma_{\rho,\sigma}/\underline\mu_\rho$
($\simeq\underline N_\sigma^{\unsure{\natural|\sharp}}
\simeq\mathbb O_{12}^{*}\cdot\mathbb G_a$).
On a cependant

\[
\underline M_{\rho,\sigma}(-\bar\jmath)\cdot\underline\mu_\rho
=\underline M_{\rho,\sigma}(-\bar\jmath)\cdot\underline\mu_G
\]

et posant

\[
(207)\quad\begin{aligned}
\underline M_{\rho,\sigma}^{\sim}(-\bar\jmath)
&\overset{\inserted{\mathrm{d\acute ef}}}{=}
\underline M_{\rho,\sigma}(-\bar\jmath)\cdot
\underline\mu_\rho/\underline\mu_\rho
\\
&=\underset{\text{\unsure{(isomorphique}}}{\operatorname{Im}}
\left(\vcenter{\hbox{\diagram{
\begin{tikzcd}[ampersand replacement=\&,column sep=small]
\underline M_{\rho,\sigma}(-\bar\jmath)\arrow[r]
 \&\underline M_{\rho,\sigma}/\mu_\rho
\end{tikzcd}
}}}\right),
\end{aligned}
\]

\inserted{$\underline\Gamma_{\rho,\sigma}^{\sim}
=\underline\Gamma_{\rho,\sigma}/\underline\mu_\rho
=(\underline M_{\rho,\sigma}/\underline\mu_\rho)/
\underline{\unsure{\mathcal G|\widetilde G}}^{+}$}

\margin{89}

\scanpage{91}
\mspage{693}

\margin{\uline{NB} On a
\diagram{\begin{tikzcd}[ampersand replacement=\&,column sep=small]
\underline\Gamma_{\rho,\sigma}^{\sim}\arrow[r]
 \&\underline N_\rho/\unsure{\underline\mu_\rho}
\end{tikzcd}}
épi de noyau $\simeq$
\[
\operatorname{Ker}\left(
\vcenter{\hbox{\diagram{\begin{tikzcd}[ampersand replacement=\&,column sep=small]
\underline N_\sigma\arrow[r]
 \&\underbrace{\underline{\mathbb O}_{12}^{*}}_{\gamma_{\rho,\sigma}}
\end{tikzcd}}}}
\right)\simeq\mathbb G_a,
\]
et On a \unsure{diagram.} \unsure{cartésien}
\diagram{\begin{tikzcd}[ampersand replacement=\&,column sep=small,row sep=large]
\underline\Gamma_{\rho,\sigma}^{\sim}\arrow[r]
 \&\underline N_\rho/\underline\mu_\rho\arrow[r]\&1
\\
\begin{gathered}
\deleted{\unsure{\underline\mu}}
\underline\Gamma_{\rho,\deleted{\unsure{\sigma}}}^{\deleted{\unsure{\sim}}}
\\[-2pt]
\simeq\underline M_{\rho,\sigma}(-\bar\jmath)
\\[-2pt]
\simeq\underline M_{\rho,\sigma}^{\sim}(-\bar\jmath)
\end{gathered}\arrow[u]\arrow[r]
 \&\underline N_\rho\arrow[u,"\mathrm{can}"]\arrow[r]\&1
\end{tikzcd}}
\unsure{le diag.} des sous-groupes de $\underline N_\rho$ du $7^\circ$
définit \unsure{un} donc \unsure{diag.} des
ss-groupes de $\underline M_{\rho,\sigma}(-\bar\jmath)$\ldots}

on aura

\[
(208)\quad\diagram{\begin{tikzcd}[ampersand replacement=\&,column sep=small]
\unsure{1}\arrow[r]\&\underline\mu_G\arrow[r]
 \&\underline M_{\rho,\sigma}^{\sim}
 (-\unsure{\bar\jmath|\deleted{\frac12}\inserted{\bar\jmath}})\arrow[r]
 \&\underline\Gamma_{\rho,\sigma}^{\sim}\deleted{\unsure{/\mu_\rho}}\arrow[r]\&1
\end{tikzcd}}
\]

i.e.
$\underline M_{\rho,\sigma}^{\sim}(-\frac12)$ est une
\uline{bisection} de
$\underline M_{\rho,\sigma}^{\sim}\overset{\mathrm{d\acute ef}}{=}
\underline M_{\rho,\sigma}/\underline\mu_\rho$
au dessus de $\underline\Gamma_{\rho,\sigma}^{\sim}$. Posant

\[
(209)\quad
\underline M_{\rho,\sigma}^{\prime\sim}(-\bar\jmath)
\overset{\mathrm{d\acute ef}}{=}
\underline M_{\rho,\sigma}^{\sim}
(-\unsure{\bar\jmath|\deleted{\frac12}\inserted{\bar\jmath}})
\cap
\underbrace{\underline M_{\rho,\sigma}^{\sim\unsure{\natural|\sharp}}}
_{\substack{\Vert\,\mathrm{d\acute ef}\\
\underline M_{\rho,\sigma}^{\unsure{\natural|\sharp}}/\mu_\rho}}
\deleted{\unsure{/\mu_\rho}}
\]

on trouve

\margin{$\underline M_{\rho,\sigma}^{\prime\sim}(-\bar\jmath)$
d'indice $4$ dans
$\underline M_{\rho,\sigma}^{\sim}(-\bar\jmath)$}

\[
(209)\quad\diagram{\begin{tikzcd}[ampersand replacement=\&,column sep=small]
\underline M_{\rho,\sigma}^{\prime\sim}(-\bar\jmath)\arrow[r,"\sim"]
 \&\underset{\substack{\Vert\,\mathrm{d\acute ef}\\
 \operatorname{Ker}\left(\underline\Gamma_{\rho,\sigma}^{\sim}
 \xrightarrow{\varepsilon_2}\underline\mu_{\unsure{S_0}}\right)}}
 {\underline\Gamma_{\rho,\sigma}^{\prime\sim}}
 \arrow[r,phantom,"\overset{\deleted{\unsure{\mathrm{d\acute ef}}}}{\subset}"]
 \&\underline\Gamma_{\rho,\sigma}^{\sim}
\end{tikzcd}},
\]

section partielle de $\underline M_{\rho,\sigma}^{\sim}$ au dessus de
$\underline\Gamma_{\rho,\sigma}^{\sim}$,
\unsure{d'im.} \unsure{p|pr}
$\Gamma_{\rho,\sigma}^{\prime\sim}$ d'indice $2$ dans
$\underline\Gamma_{\rho,\sigma}^{\prime}$.
On aura donc

\[
(210)\quad\left\{\begin{aligned}
\underline M_{\rho,\sigma}^{\sim\prime}
&\overset{\mathrm{d\acute ef}}{\unsure{=|\simeq}}
\operatorname{Ker}\left(
\underline M_{\rho,\sigma}^{\sim}\xrightarrow{\varepsilon_2}
\underline\mu_{\unsure{S_0}}
\right)
\simeq M_{\rho,\sigma}^{\prime\sim}(-\bar\jmath)\cdot
\underline{\unsure{\mathcal G|\widetilde G}}^{+}
\\
\underline M_{\rho,\sigma}^{\sim\prime\unsure{\natural|\sharp}}
&\overset{\mathrm{d\acute ef}}{=}
\operatorname{Ker}\left(
\underline M_{\rho,\sigma}^{\sim\unsure{\natural|\sharp}}
\xrightarrow{\varepsilon_2}\underline\mu
\right)
\simeq\underline M_{\rho,\sigma}^{\prime\sim}(-\bar\jmath)\cdot\underline M.
\end{aligned}\right.
\]

\uline{Remarque} Je m'arrête ici dans les développements de sous-groupes
intéressants de $\underline M_{\rho,\sigma}^{\sim}(-\bar\jmath)$
--- j'y reviendrai encore une fois, sans doute, dans le contexte des
relations aux groupes $\underline M_{\rho,\sigma}$ généraux.
Il devient clair qu'il faudra choisir,

\margin{90}

\scanpage{92}
% Main text, lines 1--7; continuation from scan 91.
par la suite, entre les notations possibles
respectives pour les groupes notés ici
$\underline M_{\rho,\sigma}$ et
$\underline M_{\rho,\sigma}/\underline\mu_\rho\simeq
\underline M_{\rho,\sigma}^{\sim}$
(et itou avec $\unsure{\natural|\sharp}$)
--- ce sont ces derniers, il est vrai, qui me semblent les plus intéressants,
car c'est eux$^{*}$ qui \unsure{rejoignent} le groupe universel
$\unsure{M|\mathcal M}=\unsure{N|\mathcal N}_{1,1}^{\sim}$
(extension de $\Gamma_{\mathbb Q}^{\sim}$ par
$\inserted{\unsure{\mathcal G|G}_{0}^{\wedge}=}
(\mathrm{SL}(2,\mathbb Z))^{\wedge}$).
\margin{plutôt, le groupe de leurs $\widehat{\mathbb Z}$-points
\unsure{respct...|respect...}}

% Main text, lines 8--25.
Il semblerait que ces groupes \unsure{mériteraient} la notation la plus simple,
en l'occurrence $\underline M_{\rho,\sigma}$, et les groupes notés
précédemment $\underline M_{\rho,\sigma}$ prendraient leur $\sim$,
ou un signe analogue. (Le $\sim$ est déjà utilisé dans un sens assez
différent dans $\Gamma_{\mathbb Q}^{\sim}$,
$M_{0,3}^{\sim}$ etc --- d'autre part les exposants
$',\deleted{\unsure{\natural|\sharp}},'',\,''',\,{}^{\circ},\,+,\,!,\,
\unsure{\natural|\sharp}$
\deleted{$\unsure{',\,',\,',\,',\,+}$}
sont déjà utilisés dans des sens différents.
Le double emploi du $\sim$ risque d'ailleurs d'entraîner des conflits de
notation tôt ou tard, il est peut-être préférable de remplacer le $\sim$
dans $\Gamma_{\mathbb Q}^{\sim}$ par des lettres telles que
$\Gamma_{\mathbb Q}^{\text{ét}}$
(suggérant ``étendu'') ou \deleted{\unsure{gén}}
``généralisé'' \unsure{gén\ldots}\ \deleted{\illegible{1}}
Il faudra que je résolve ces di-
\margin{91}

\scanpage{93}
\mspage{694}
% Main text, lines 1--5; continuation of di- on scan 92.
lemmes notationnels, pour l'instant je garde les notations introduites pour
\unsure{le besoin|les besoins} dans la présente section, qui seront
\unsure{revisées|révisées} avec soin dans un \S\ spécial ultérieur.

% Main text, lines 6--9.
Il faudrait encore décrire les sous-groupes
$\underline M_{\rho,\sigma}^{\sim}(-\bar j)$,
$\underline M_{\rho,\sigma}^{\prime\sim}(-\bar j)$
dans $\underline M_{\rho,\sigma}^{\sim}$ écrit comme produit
$\frac12$ direct, on trouve
\[
\begin{aligned}
(211)\qquad
\underline M_{\rho,\sigma}^{\sim}(-\bar j)
&\deleted{\unsure{\simeq\underline{\mathbb O}_{12}^{*}
(\mathbb G_a\times\underline\mu_G)}}
\inserted{\deleted{\unsure{\underline\mu(\lambda,x)}}}
\subset \underline M_{\rho,\sigma}^{\sim}
\simeq \mathbb O_{12}^{*}
(\mathbb G_a\times\underline{\unsure{\mathcal G|\widetilde G}}^{+})
\\[-1pt]
&=\left\{
\deleted{\underset{\mathbb O_{12}^{*}}
{\unsure{\underline\mu}}}
\left(
\underset{\substack{\rotatebox{90}{$\in$}\\\mathbb G_a}}{\lambda},
\underset{\substack{\rotatebox{90}{$\in$}\\
\underline{\unsure{\mathcal G|\widetilde G}}^{+}}}{x}
\right)\cdot
\underset{\substack{\rotatebox{90}{$\in$}\\
\mathbb O_{12}^{*}}}{\underline\mu}
\ \middle|\ x\cdot\deleted{\unsure{\underline\mu}}
\underbrace{u_{\underline\mu}\deleted{\unsure{\rho}}}
_{\in\underline G_{\rho,\sigma}}
\in N_\rho^{?}\right\}
\end{aligned}
\]
% Main text and calculation, lines 12--15.
où $u_{\underline\mu}$ est défini dans (151), en particulier
\[
u_{\underline\mu}(\rho)=\lambda_1^{\nu}
\unsure{\{\rho\}|\rho}\deleted{\unsure{=}}
\qquad\text{où}\quad \nu=\nu(\underline\mu),
\]
\[
\deleted{=\lambda_1^{3\nu''}\lambda_1^{-\nu_3}\rho}
\]
\deleted{donc}
\diagram{\begin{tikzcd}[ampersand replacement=\&,column sep=0.1em,row sep=1.6em,nodes in empty cells]
\deleted{(x} \&
\deleted{u_{\underline\mu})(\rho)=x(\lambda_1)^{\unsure{\nu}}x(\rho)}
\\
\inserted{\unsure{\delta^{\nu''}\tau^{\nu_3}}}
\arrow[u,bend left=20] \&
\end{tikzcd}}
% Replacement calculation, lines 15--18.
on a vu (154)
\[
v=\underbrace{\sigma^{\nu_3}\delta^{-\nu''}u_{\underline\mu}}
_{\unsure{v|u}}\in\Gamma\underline N_\rho,
\qquad \underline\chi(v)=\underline\mu
\]
\margin{où $\delta=\lambda_1\lambda_0^{-1}$.}
et posant $\sigma=\omega\tau'\tau$, le premier
\unsure{$\mathrm m^{\mathrm b}$|mb} s'écrit aussi
\[
\deleted{(\unsure{\omega\tau'})^{\nu_3}\tau^{\nu_3}\unsure{u}}
\]
\[
(212)\qquad
v={\tau'}^{\nu_3}\left(\omega^{\nu_3}\cdot
\underbrace{\tau^{\nu_3}\delta^{-\nu''}u_{\underline\mu}}
_{\substack{
\in\Gamma\underline{\deleted{\unsure{N}}\inserted{G}}
_{\rho,\sigma}^{\unsure{\natural|\sharp}},\ \text{car les}\\
\text{trois facteurs le}\\
\text{sont; donc}\\
\in\Gamma\underline N_\rho\cap
\underline G_{\rho,\sigma}^{\unsure{\natural|\sharp}}
=\Gamma\underline Z_\rho^{\unsure{\natural|\sharp}}
}}\right)
\]
\margin{92}

\scanpage{94}
% Main text, lines 1--4.
et on sait que $v\in\deleted{\unsure{\underline Z}}
N_\rho^{?}$, donc la forme
\deleted{les \unsure{éléments} de}
générale des \deleted{éléments}\inserted{sections}
de $N_\rho^{?}$ au dessus de
$\underline\mu=(\nu_2,\nu_3,\unsure{\gamma_0|\nu_0})
\in\Gamma\mathbb O_{12}^{*}$
\[
(212)\qquad
(\underline N_\rho^{?})_{\underline\mu}
=\left\{\omega^{i}\unsure{\sigma}^{\nu_3}\delta^{-\nu''}
u_{\underline\mu}\ \middle|\
i\in\Gamma\{0,1\}_{S_0}\right\}
\]
\unsure{avec|où}
$\nu_3=\nu_3(\underline\mu)$,
$\nu''=\nu''(\underline\mu)$,
$\delta=\lambda_1\lambda_0^{-1}$. Donc on trouve
\[
\begin{aligned}
(213)\qquad
\underline M_{\rho,\sigma}^{\sim}(-\bar j)
=\left\{\left(
\underset{\substack{\rotatebox{90}{$\in$}\\\mathbb G_a}}{\lambda},
\underset{\substack{\rotatebox{90}{$\in$}\\
\underline{\unsure{\mathcal G|\widetilde G}}^{+}}}{x}
\right)\cdot
\underset{\substack{\rotatebox{90}{$\in$}\\
\mathbb O_{12}^{*}}}{\underline\mu}
\ \middle|\
\begin{gathered}
x=\omega^{i}\sigma^{\nu_3}\delta^{-\nu''}u_{\underline\mu},\ \text{où}\\
i\in\Gamma\{0,1\}_{S_0}
\end{gathered}\right\}
\end{aligned}
\]
\margin{$\delta=\lambda_1\lambda_0^{-1}$}

On aura donc
\[
\begin{aligned}
(214)\qquad
\underline M_{\rho,\sigma}^{\prime\sim}(-\bar j)
&=\underline M_{\rho,\sigma}^{\sim}(-\bar j)\cap
\underline M_{\rho,\sigma}^{\unsure{\natural|\sharp}}\\
&=\left\{(\lambda,\delta^{-\nu''})\underline\mu
\ \middle|\
\begin{gathered}
\underline\mu\in\Gamma\mathbb O_{12}^{*\prime}
\quad(\text{donc }\varepsilon_3(\underline\mu)=\unsure{0|1})\\
\lambda\in\Gamma\mathbb G_a\\
\nu''=\nu''(\underline\mu)
\end{gathered}\right\}
\end{aligned}
\]
\margin{93}

\scanpage{95}
\mspage{695}
% Heading and opening, lines 1--3.
$11^{\mathrm o})$ L'homomorphisme
$M\longrightarrow\underline M_{\rho,\sigma}(\widehat{\mathbb Z})^{\Delta}$,
relation avec la tour des courbes de Fermat.

Notons qu'on a pour tout entier $N\in\mathbb N^{*}$
\[
(215)\qquad
\mathbb O_{N}^{\Delta}
(\widehat{\mathbb Z})\simeq\widehat{\mathbb Z}
\qquad\text{d'où}\qquad
\mathbb O_{N}^{*}(\widehat{\mathbb Z})^{\Delta}
\simeq\widehat{\mathbb Z}^{*}.
\]
\margin{Le $\Delta$ indique qu'on prend les sections qui sont constantes
mod comp. neutre.}

On a donc, compte tenu des décompositions en produits semi-directs (150) etc,
des isomorphismes
\[
(216)\qquad
\left\{\begin{aligned}
Z_{\rho,\sigma}^{\unsure{\ell|\wedge}}
&\overset{\mathrm{dfn}}{=}
\underline Z_{\rho,\sigma}^{\deleted{\unsure{\natural|\sharp}}}
(\widehat{\mathbb Z})^{\Delta}
=\widehat{\mathbb Z}^{*}\cdot
\deleted{\unsure{\widehat{\mathbb Z}}}\widehat L_0
\quad\text{où}\quad
\widehat L_0\simeq\widehat{\mathbb Z}\cdot\lambda_0\simeq\widehat{\mathbb Z}
\\
G_{\rho,\sigma}
&\overset{\mathrm{dfn}}{=}
\underline G_{\rho,\sigma}^{\deleted{\unsure{\Delta|\natural}}}
(\widehat{\mathbb Z})^{\Delta}
=\widehat{\mathbb Z}^{*}\cdot
\underline{\unsure{\mathcal G|\widetilde G}}^{+\wedge}
\\
G_{\rho,\sigma}^{\unsure{\natural|\sharp}}
&\overset{\mathrm{dfn}}{=}
\underline G_{\rho,\sigma}^{\unsure{\natural|\sharp}}
(\widehat{\mathbb Z})^{\Delta}
\simeq\widehat{\mathbb Z}^{*}\cdot\widehat M
\\
N_{\rho\deleted{\unsure{,\sigma}}}^{\sim\deleted{\unsure{\natural}}}
&\overset{\mathrm{df}}{=}
\underline N_\rho^{\sim}(\widehat{\mathbb Z})^{\Delta}
\simeq\mathbb O_{12}^{*}
(\widehat{\mathbb Z})^{\Delta}\simeq\widehat{\mathbb Z}^{*}
\\
N_\sigma^{\unsure{\natural|\sharp}}
&\overset{\mathrm{df}}{=}
\underline N_\sigma^{\unsure{\natural|\sharp}}
(\widehat{\mathbb Z})^{\Delta}
\simeq\deleted{\mathbb O_{12}^{*}}
\widehat{\mathbb Z}^{*}\cdot\widehat{\mathbb Z}
\\
\underset{\deleted{\unsure{\infty}}}{M}_{\rho,\sigma}^{\sim}
&\overset{\mathrm{df}}{=}
\underline M_{\rho,\sigma}^{\deleted{\unsure{\natural|\sharp}}}
(\widehat{\mathbb Z})^{\Delta}
\simeq\widehat{\mathbb Z}^{*}\cdot
(\unsure{\widehat{\mathbb Z}}\times
\unsure{\mathcal G|\widetilde G}^{+\wedge})
\\
M_{\rho,\sigma}^{\sim\unsure{\natural|\sharp}}
&\overset{\mathrm{df}}{=}
\underline M_{\rho,\sigma}^{\unsure{\natural|\sharp}}
(\widehat{\mathbb Z})^{\Delta}
\simeq\widehat{\mathbb Z}^{*}\cdot
(\widehat{\mathbb Z}\times\widehat M)
\end{aligned}\right.
\]
% Main text, lines below (216).
où les opérations de $\widehat{\mathbb Z}^{*}$ sur les
\unsure{$2^{\text{ièmes}}$|$2^{\text{èmes}}$} facteurs
$\frac12$ directs sont explicitées dans les sections précédentes.
On notera que, malgré les \unsure{apprêtages|apparences} des
notations et des définitions, les groupes \inserted{profinis}
$Z_{\rho,\sigma}$ et précédents sont aussi ceux qui seraient
\unsure{construits} \deleted{\unsure{$\mathcal G$}}
directement, en termes de
$(\unsure{\mathcal G|\widetilde G}^{\unsure{+|\natural}\wedge},
\unsure{\mathcal G|\widetilde G}^{\unsure{\natural|\sharp}\wedge}
=M^{\wedge})$
(munis$^{*}$ de $\rho,\sigma,\tau$), dans \S\ 49 $\uline{\overline{\mathrm{VI}}}$.
On a donc un hom. canonique
\diagram{\[
(217)\qquad
\left\{\begin{gathered}
\begin{tikzcd}[ampersand replacement=\&,column sep=5em,row sep=2.5em]
M=\unsure{N|\mathcal N}_{1,1}^{\text{ét}}
\arrow[r]\arrow[d,phantom,"\cup" description]
\& M_{\rho,\sigma}^{\sim}\arrow[d,phantom,"\cup" description]
\\
M^{!}\arrow[r] \& M_{\rho,\sigma}^{\sim\unsure{\natural|\sharp}}
\end{tikzcd}
\end{gathered}\right.
\]}
qui comme d'habitude se décompose en
\margin{94}
